%! Displaying and Describing Distributions — Unit 2, Lesson 2.1 — Student Packet Cover Sheet
% THREE packet rows — Warm-Up, Guided Notes & Practice, Homework. The homework is the individual
% practice, assigned at Close & Assign and done outside class (Unit 2 on), and it is SCORED, so
% its cell is a \blank{}, never NA. Standard codes stay in the plan, never here.
\documentclass[12pt]{article}
\usepackage{saar-article}
\usepackage{saar-boxes}
\usepackage{ltablex}
\keepXColumns
% Measured banner: the band grows to fit the title block, so a title long enough to wrap grows
% the band rather than running out of it (the stats course's \coverbanner; saar-article does not
% carry one, so it is defined here — every cover copies this block verbatim).
\newsavebox{\bannerbox}
\newlength{\bannerht}
\newlength{\coverbannerpad}
\setlength{\coverbannerpad}{0.16in}       % breathing room above and below the text
\newcommand{\coverbanner}[2]{%
  \sbox{\bannerbox}{%
    \begin{minipage}{\dimexpr\paperwidth-1.4in\relax}%
      \centering
      {\color{white}\textbf{\LARGE \CourseName}}\\[4pt]
      {\color{frostmid}\large\textbf{#1}}\\[6pt]
      {\color{white}\textbf{\large #2}}%
    \end{minipage}}%
  \setlength{\bannerht}{\dimexpr\ht\bannerbox+\dp\bannerbox+2\coverbannerpad\relax}%
  \noindent\begin{tikzpicture}[remember picture, overlay]
    \fill[cerulean] (current page.north west)
      rectangle ([yshift=-\bannerht]current page.north east);
    \node[anchor=north, inner sep=0pt]
      at ([yshift=-\coverbannerpad]current page.north) {\usebox{\bannerbox}};
  \end{tikzpicture}\par
  % The text body starts 0.6in below the page edge (geometry top=0.6in), so reserve only what
  % the band overruns that, plus a gap under the band.
  \vspace*{\dimexpr\bannerht-0.6in+0.16in\relax}%
}

\begin{document}

% Full-bleed cerulean banner, sized to the title block.
\coverbanner{Unit 2}{Lesson 2.1 \quad Displaying and Describing Distributions}

\vspace{0.06in}
% NAMESTRIP: this is the ONLY component that carries the name row.
\namedateperiod
\vspace{0.18in}

% GRADUAL RELEASE: the vocabulary is taught in the guided notes, so the targets name it
% plainly. The standard codes (PS.DS.1a-b) stay in the teacher-facing plan.
\begin{learningtargetbox}
\small
\begin{enumerate}[leftmargin=1.5em, itemsep=5pt]
  \item \ldots{} read a \textbf{dot plot}, a \textbf{stemplot} (using its \textbf{key}), and a
        \textbf{histogram} from a calculator or Desmos, and say what one dot, leaf, or bar stands
        for.
  \item \ldots{} describe a \textbf{distribution} by its \textbf{shape} (\textbf{symmetric},
        \textbf{skewed}, \textbf{uniform}; \textbf{unimodal} or \textbf{bimodal}), its
        \textbf{center}, and its \textbf{spread} (\textbf{range}), in a sentence with context and
        units.
  \item \ldots{} point out unusual features --- \textbf{outliers}, \textbf{gaps}, and
        \textbf{clusters}.
  \item \ldots{} name the direction of \textbf{skew} by the \textbf{tail}, not by the peak.
\end{enumerate}
\end{learningtargetbox}

\vspace{0.14in}

\begin{tocbox}
\small
\renewcommand{\arraystretch}{1.5}
\begin{tabularx}{\linewidth}{@{} c l X r @{}}
\rowcolor{frost}
\textbf{\#} & \textbf{Component} & \textbf{Description} & \textbf{Score} \\
\midrule
1 & Warm-Up & Back to Unit 1: quantitative or categorical, a variable's units, and a
             frequency table of trip times & \blank{1.2cm} \\
2 & Guided Notes \& Practice & Dot plots, stemplots, and histograms; describing a
             distribution by shape, unusual features, center, and spread; which way a skew
             points --- then one team's free throws we describe together & \blank{1.2cm} \\
% Row 3 is the lesson's GRADED practice, assigned at the close and done outside class. Packet
% pages are the default; a Desmos activity may replace them on a given day, decided at Close &
% Assign. The row and its score cell never change.
3 & Homework & Oak Hollow Library summer reading: finish a stemplot, read a histogram, and
             name and explain two skews ---
             \textbf{scored, assigned today, due the first class after two study halls} & \blank{1.2cm} \\
\midrule
  & \multicolumn{2}{r}{\textbf{Total}} & \blank{1.2cm} \\
\end{tabularx}
\end{tocbox}

\vspace{0.14in}

% Keep in Mind carries the lesson's CONTENT — the rule, the test, the trap — never its process.
\begin{remindbox}
\small
To describe a \textbf{distribution}, name its \textbf{shape}, any \textbf{unusual features}
(outliers, gaps, clusters), its \textbf{center}, and its \textbf{spread} --- always with the
individuals, the variable, and the units. \textbf{Skew is named by the tail:} if the long tail
stretches toward the larger values, the distribution is \textbf{skewed right}; toward the smaller
values, \textbf{skewed left}. \textbf{Watch out:} the peak --- where most of the data pile up ---
sits on the \emph{opposite} side from the tail. A pile on the left with a tail to the right is
skewed \emph{right}.
\end{remindbox}

\end{document}
